\documentclass[11pt,a4paper]{article}

\usepackage[utf8]{inputenc}
\usepackage[T1]{fontenc}
\usepackage{amsmath,amssymb,amsfonts}
\usepackage{graphicx}
\usepackage{booktabs}
\usepackage{hyperref}
\usepackage{geometry}
\usepackage{algorithm}
\usepackage{algpseudocode}
\usepackage{xcolor}
\usepackage{caption}
\usepackage{subcaption}
\usepackage{microtype}
\usepackage{bm}
\usepackage{mathtools}
\usepackage{cleveref}

\geometry{margin=2.5cm}

\hypersetup{
    colorlinks=true,
    linkcolor=blue!60!black,
    citecolor=green!50!black,
    urlcolor=blue!60!black
}

\title{%
  \Large \textbf{Multiple Instance Learning on CT Patch Embeddings\\
  for Lung Cancer Survival Prediction}\\[0.5em]
  \large Technical Report
}
\author{}
\date{\today}

\begin{document}

\maketitle
\tableofcontents
\newpage

% ============================================================
\section{Introduction}
% ============================================================

Predicting overall survival (OS) from computed tomography (CT) imaging is a central
challenge in thoracic oncology. For patients with non-small-cell lung cancer (NSCLC)
or small-cell lung cancer (SCLC), CT scans encode rich morphological and textural
information about the primary tumour and its surrounding tissue that is prognostically
relevant. Classical radiomics pipelines extract hand-crafted features from
segmentation masks and feed them into shallow classifiers; however, they discard spatial
context outside the annotated region of interest and are sensitive to segmentation
quality.

The framework described in this report takes a different approach: it decomposes
each CT volume into a set of three-dimensional (3-D) sub-volume patches, encodes
every patch with a pretrained CT foundation model, and then aggregates the resulting
bag of embeddings at the patient level using a \emph{Multiple Instance Learning}
(MIL) classifier. This design has three key advantages:
\begin{enumerate}
  \item No pixel-level annotation is required; supervision is provided only at the
        patient level (the binary survival label).
  \item Whole-volume context is preserved, allowing the model to integrate information
        from the tumour core, margins, and surrounding parenchyma simultaneously.
  \item The modular architecture supports multiple aggregation strategies that can be
        compared systematically under a common training and evaluation pipeline.
\end{enumerate}

The primary model highlighted in this report is the \textbf{LSTM-based MIL} variant.
Patches are ordered anatomically along the superior-inferior (Z) axis before being
fed as a temporal sequence to a bidirectional Long Short-Term Memory (LSTM) network.
The LSTM hidden state is conceptualised as a \emph{risk accumulator}: as the network
scans the volume from apex to base, its memory progressively integrates local patch-level
lesion signals into a coherent, holistic representation of cumulative tumour burden.
This narrative motivates the architectural choices described in
\cref{sec:lstm}.

% ============================================================
\section{Problem Formulation}
\label{sec:problem}
% ============================================================

\subsection{Multiple Instance Learning}

Let $\mathcal{D} = \{(B_i, y_i)\}_{i=1}^{N}$ denote a dataset of $N$ patients,
where $y_i \in \{0, 1\}$ is the binary patient-level label and
$B_i = \{\mathbf{x}_{i,1}, \ldots, \mathbf{x}_{i,K_i}\}$ is the corresponding
\emph{bag} of $K_i$ patch embeddings, with
$\mathbf{x}_{i,k} \in \mathbb{R}^{d}$.

Under the standard MIL assumption, the bag label is observed but instance
labels are latent. The primary MIL assumption used here is the \emph{collective}
(or embedding-space) variant~\cite{carbonneau2018}: the bag-level representation
is obtained by aggregating instance embeddings, and the classifier operates on this
aggregate:
\begin{equation}
  \hat{y}_i = f_\theta\!\left(\phi\!\left(\{\mathbf{x}_{i,k}\}_{k=1}^{K_i}\right)\right),
\end{equation}
where $\phi(\cdot)$ is the aggregation operator (LSTM, Transformer, pooling, etc.)
and $f_\theta(\cdot)$ is a learned classifier parameterised by $\theta$.

\subsection{Survival Prediction Endpoints}

Two binary classification tasks are defined:
\begin{align}
  y^{\text{OS6}}_i  &= \mathbf{1}[\text{OS}_i \leq 6  \text{ months}], \\
  y^{\text{OS24}}_i &= \mathbf{1}[\text{OS}_i \leq 24 \text{ months}],
\end{align}
where $\text{OS}_i$ is the overall survival time in months for patient $i$.
These thresholds define clinically meaningful risk strata: OS6 separates patients
with very poor short-term prognosis, while OS24 separates those who achieve
longer-term disease control.

% ============================================================
\section{Dataset and Preprocessing}
\label{sec:data}
% ============================================================

\subsection{Patient Cohort}

The dataset consists of CT volumes from patients diagnosed with NSCLC or SCLC,
collected from \emph{multiple clinical centres}. Each sample is represented as a
dictionary containing:
\begin{itemize}
  \item \texttt{features}: a tensor of shape $[K, d]$ where $K \leq K_{\max}$ is the
        number of patches and $d = 512$ is the embedding dimensionality;
  \item \texttt{OS\_6} / \texttt{OS\_24}: the binary survival label;
  \item \texttt{patient\_id}: a unique patient identifier;
  \item \texttt{center}: the originating clinical centre, used for subgroup evaluation.
\end{itemize}

\subsection{Class Imbalance}

The dataset exhibits mild class imbalance between the positive (short-survivor) and
negative (long-survivor) classes. To mitigate this, an \emph{oversampling} strategy
is applied to the minority class during training. The oversampling factor
$\alpha \geq 1$ controls how many times each minority-class sample is replicated,
yielding an effective training set size of:
\begin{equation}
  N_{\text{eff}} = N_{\text{maj}} + \alpha \cdot N_{\text{min}},
\end{equation}
where $N_{\text{maj}}$ and $N_{\text{min}}$ are the majority and minority class
counts respectively. Additionally, the cross-entropy loss is optionally weighted
by inverse class frequencies.

\subsection{Dataset Splits}

Data are split into training, validation, and test partitions. Two modes are
supported:
\begin{itemize}
  \item \textbf{Random split}: a proportion $r_{\text{val}}$ (default 0.15) is held out
        for validation and $r_{\text{test}}$ (default 0.15) for testing, with
        stratification by label.
  \item \textbf{Pre-split}: fixed \texttt{train\_set.pkl}, \texttt{val\_set.pkl},
        and \texttt{test\_set.pkl} files are provided, enabling reproducible
        multi-centre evaluation.
\end{itemize}

\subsection{Variable-Length Bags and Padding}

Because patches are extracted from the actual extent of each CT volume, the bag
size $K_i$ varies across patients. Bags are zero-padded to a common length
$K_{\max} = 300$ for batched processing:
\begin{equation}
  \tilde{\mathbf{X}}_i \in \mathbb{R}^{K_{\max} \times d},
  \quad
  \tilde{\mathbf{X}}_i[k] =
  \begin{cases}
    \mathbf{x}_{i,k} & k \leq K_i, \\
    \mathbf{0}        & k > K_i.
  \end{cases}
\end{equation}
A binary padding mask $\mathbf{m}_i \in \{0,1\}^{K_{\max}}$ is derived
post-projection by testing whether the $\ell^1$-norm of each token is non-zero:
\begin{equation}
  m_{i,k} = \mathbf{1}\!\left[\|\tilde{\mathbf{x}}_{i,k}\|_1 > 0\right].
\end{equation}
The mask is used to suppress padded positions in subsequent layers.

% ============================================================
\section{Patch Embedding Extraction}
\label{sec:embeddings}
% ============================================================

\subsection{CT-FM Foundation Model}

Each CT volume is partitioned into non-overlapping (or mildly overlapping)
three-dimensional sub-volume crops. Every crop $\mathbf{v}_{i,k} \in \mathbb{R}^{H \times W \times D}$
is encoded by \textbf{CT-FM}~\cite{ctfm}, a vision foundation model pretrained
in a self-supervised manner on over 140{,}000 CT scans spanning all anatomical sites.
CT-FM produces a fixed-dimensional embedding for each crop:
\begin{equation}
  \mathbf{x}_{i,k} = \text{CT-FM}(\mathbf{v}_{i,k}) \in \mathbb{R}^{512}.
\end{equation}
The encoder weights are \emph{frozen} during MIL training; the framework operates
entirely in the embedding space produced by CT-FM.

\subsection{Assumptions on the Embedding Space}

\begin{enumerate}
  \item \textbf{Semantic sufficiency}: CT-FM embeddings capture radiologically
        meaningful variation (texture, morphology, attenuation patterns) at the local
        patch level, sufficient for survival-relevant discrimination.
  \item \textbf{Domain alignment}: despite being pretrained on a heterogeneous
        multi-anatomy corpus, CT-FM features transfer effectively to the thoracic oncology domain.
  \item \textbf{Patch independence} (at the encoder stage): each patch is encoded
        independently; spatial context across patches is modelled exclusively by the
        MIL aggregator.
\end{enumerate}

\subsection{Anatomical Ordering}

Before being passed to the LSTM aggregator, patches within each bag are sorted
by their anatomical position along the superior-inferior (Z) axis of the CT volume.
This imposes a clinically grounded sequential structure: the model reads the thorax
from apex to base, mirroring the typical radiologist inspection order and enabling
the LSTM to exploit spatial continuity.

% ============================================================
\section{Model Architectures}
\label{sec:models}
% ============================================================

All models share the same input interface: a padded bag tensor
$\tilde{\mathbf{X}} \in \mathbb{R}^{B \times K_{\max} \times d}$ where $B$ is the
batch size, and produce logits $\mathbf{z} \in \mathbb{R}^{B \times C}$ where
$C = 2$ is the number of classes.

% ------------------------------------------------------------
\subsection{LSTM-Based MIL (Primary Model)}
\label{sec:lstm}
% ------------------------------------------------------------

\subsubsection{Motivation}

The LSTM-based MIL model treats the ordered sequence of patch embeddings as a
\emph{temporal signal} to be processed recurrently. As the network scans the
volume from apex to base, the LSTM hidden state acts as a \textbf{running risk
accumulator}: each new patch updates the memory in proportion to the local lesion
signal it carries, and the final (or attention-weighted) hidden state represents
the \emph{cumulative tumour burden} integrated over the entire volume. This
inductive bias is well-matched to the spatial continuity of lung tumours, which
typically manifest as contiguous lesions whose full extent spans multiple
anatomical levels.

\subsubsection{Feature Projection}

The input embeddings are first projected into a lower-dimensional model space:
\begin{equation}
  \mathbf{H}^{(0)} = \tilde{\mathbf{X}} \,\mathbf{W}_{\text{proj}} + \mathbf{b}_{\text{proj}},
  \quad
  \mathbf{W}_{\text{proj}} \in \mathbb{R}^{d \times d_h},
\end{equation}
yielding $\mathbf{H}^{(0)} \in \mathbb{R}^{B \times K_{\max} \times d_h}$ where
$d_h$ is the hidden dimension (default $d_h = 128$).

\subsubsection{Bidirectional LSTM}

The projected sequence is processed by an $L$-layer bidirectional LSTM
(default $L = 2$). At each layer $\ell$ and each time step $k$, the LSTM
cell computes:
\begin{align}
  \mathbf{f}_k^{(\ell)} &= \sigma\!\left(\mathbf{W}_f \mathbf{h}_{k-1}^{(\ell)} + \mathbf{U}_f \mathbf{h}_k^{(\ell-1)} + \mathbf{b}_f\right), \label{eq:forget}\\
  \mathbf{i}_k^{(\ell)} &= \sigma\!\left(\mathbf{W}_i \mathbf{h}_{k-1}^{(\ell)} + \mathbf{U}_i \mathbf{h}_k^{(\ell-1)} + \mathbf{b}_i\right), \label{eq:input}\\
  \mathbf{o}_k^{(\ell)} &= \sigma\!\left(\mathbf{W}_o \mathbf{h}_{k-1}^{(\ell)} + \mathbf{U}_o \mathbf{h}_k^{(\ell-1)} + \mathbf{b}_o\right), \label{eq:output}\\
  \tilde{\mathbf{c}}_k^{(\ell)} &= \tanh\!\left(\mathbf{W}_c \mathbf{h}_{k-1}^{(\ell)} + \mathbf{U}_c \mathbf{h}_k^{(\ell-1)} + \mathbf{b}_c\right), \label{eq:cell_cand}\\
  \mathbf{c}_k^{(\ell)} &= \mathbf{f}_k^{(\ell)} \odot \mathbf{c}_{k-1}^{(\ell)} + \mathbf{i}_k^{(\ell)} \odot \tilde{\mathbf{c}}_k^{(\ell)}, \label{eq:cell}\\
  \mathbf{h}_k^{(\ell)} &= \mathbf{o}_k^{(\ell)} \odot \tanh\!\left(\mathbf{c}_k^{(\ell)}\right), \label{eq:hidden}
\end{align}
where $\sigma(\cdot)$ is the sigmoid function, $\odot$ denotes element-wise
multiplication, and the superscript $(\ell)$ indexes the layer.

The gates in \cref{eq:forget,eq:input,eq:output,eq:cell_cand,eq:cell,eq:hidden}
embody the \emph{risk accumulation} interpretation:
\begin{itemize}
  \item The \textbf{forget gate} $\mathbf{f}_k$ decides how much of the accumulated
        risk memory is preserved when transitioning from patch $k-1$ to patch $k$.
        Anatomically uninformative patches (e.g., air, background) tend to induce
        $\mathbf{f}_k \approx \mathbf{1}$, preserving prior risk context unchanged.
  \item The \textbf{input gate} $\mathbf{i}_k$ controls the fraction of the new
        patch's risk signal $\tilde{\mathbf{c}}_k$ that is written into memory,
        effectively weighting the local pathological signal of each slice.
  \item The \textbf{cell state} $\mathbf{c}_k$ is the risk accumulator: it integrates
        patch-level signals across the full longitudinal extent of the volume, allowing
        the model to capture multi-level lesion distributions.
\end{itemize}

Because the LSTM is \textbf{bidirectional}, two LSTM chains are run:
one in the anatomical order (apex $\to$ base) and one in reverse
(base $\to$ apex). Their outputs are concatenated at each time step:
\begin{equation}
  \mathbf{h}_k = \left[\overrightarrow{\mathbf{h}}_k \,\Big\|\, \overleftarrow{\mathbf{h}}_k\right]
  \in \mathbb{R}^{2 d_h},
\end{equation}
where $\|\cdot\|$ denotes concatenation. This yields the full sequence of
hidden states $\mathbf{H} = [\mathbf{h}_1, \ldots, \mathbf{h}_{K_{\max}}]
\in \mathbb{R}^{B \times K_{\max} \times 2d_h}$.

\subsubsection{Padding Masking and Layer Normalisation}

The padding mask is applied to suppress the hidden states of padded positions:
\begin{equation}
  \hat{\mathbf{h}}_k = m_k \cdot \mathbf{h}_k,
\end{equation}
ensuring that zero-padded inputs do not contaminate the aggregated representation.
Layer normalisation is then applied across the feature dimension:
\begin{equation}
  \bar{\mathbf{h}}_k = \text{LayerNorm}(\hat{\mathbf{h}}_k).
\end{equation}

\subsubsection{Attention-Based Pooling}

To aggregate the sequence of normalised hidden states into a single bag-level
vector, an additive attention mechanism is applied. For each token
$\bar{\mathbf{h}}_k$, an unnormalised attention score is computed:
\begin{equation}
  u_k = \mathbf{w}_2^\top \tanh\!\left(\mathbf{W}_1 \bar{\mathbf{h}}_k + \mathbf{b}_1\right),
  \quad
  \mathbf{W}_1 \in \mathbb{R}^{d_h \times 2d_h},\;
  \mathbf{w}_2 \in \mathbb{R}^{d_h},
\end{equation}
and then normalised via softmax over the $K_{\max}$ positions:
\begin{equation}
  \alpha_k = \frac{\exp(u_k)}{\sum_{j=1}^{K_{\max}} \exp(u_j)},
  \quad k = 1, \ldots, K_{\max}.
\end{equation}
The bag-level context vector is the attention-weighted sum of hidden states:
\begin{equation}
  \mathbf{z}_{\text{bag}} = \sum_{k=1}^{K_{\max}} \alpha_k \,\bar{\mathbf{h}}_k \;\in\; \mathbb{R}^{2d_h}.
\end{equation}
The weights $\{\alpha_k\}$ are interpretable as patch-level importance scores;
they indicate which anatomical sub-volumes contributed most to the final
prediction and can be used for spatial localisation of risk-relevant regions.

When the attention module is disabled, global average pooling over non-padded
positions is used instead:
\begin{equation}
  \mathbf{z}_{\text{bag}} = \frac{1}{\sum_k m_k} \sum_{k=1}^{K_{\max}} m_k \,\bar{\mathbf{h}}_k.
\end{equation}

\subsubsection{Classifier Head}

The bag-level vector is passed through a two-layer MLP with intermediate
LayerNorm and GELU activation:
\begin{align}
  \mathbf{r} &= \text{Dropout}\!\left(\text{GELU}\!\left(\text{LayerNorm}\!\left(\mathbf{W}_3 \,\mathbf{z}_{\text{bag}} + \mathbf{b}_3\right)\right)\right),
  \quad \mathbf{W}_3 \in \mathbb{R}^{d_h \times 2d_h}, \\
  \mathbf{z} &= \mathbf{W}_4 \,\mathbf{r} + \mathbf{b}_4,
  \quad \mathbf{W}_4 \in \mathbb{R}^{C \times d_h},
\end{align}
producing unnormalised logits $\mathbf{z} \in \mathbb{R}^C$.
A final dropout is applied to $\mathbf{z}_{\text{bag}}$ before the first linear
layer to regularise the bag-level representation.

\subsubsection{Full Forward Pass Summary}

\begin{equation}
\boxed{
  \hat{y} = \arg\max_c \left[\mathbf{W}_4\,\text{GELU}\!\left(\text{LN}\!\left(\mathbf{W}_3 \sum_k \alpha_k\,\text{LN}\!\left(\text{BiLSTM}\!\left(\mathbf{W}_{\text{proj}}\,\tilde{\mathbf{X}}\right)\right)_k\right)\right)\right]_c
}
\end{equation}

\subsubsection{Weight Initialisation}

Linear layers (excluding LSTM internals) are initialised with Kaiming normal
initialisation~\cite{he2015}:
\begin{equation}
  W_{ij} \sim \mathcal{N}\!\left(0,\, \frac{2}{n_{\text{in}}}\right),
\end{equation}
where $n_{\text{in}}$ is the fan-in of the layer. Bias parameters are
initialised to zero. LSTM weights are left at their PyTorch defaults
(uniform initialisation scaled by $1/\sqrt{n_{\text{hidden}}}$).

% ------------------------------------------------------------
\subsection{Transformer-Based MIL}
\label{sec:transformer}
% ------------------------------------------------------------

The Transformer variant treats the patch sequence as a \emph{set} (order-invariant
by design, but equipped with learned positional encodings to optionally capture
spatial structure).

\paragraph{Feature Projection and Positional Encoding.}
Input embeddings are projected and augmented with a learned positional encoding
$\mathbf{P} \in \mathbb{R}^{K_{\max} \times d_h}$:
\begin{equation}
  \mathbf{H}^{(0)}_k = \mathbf{W}_{\text{proj}}\,\mathbf{x}_k + \mathbf{P}_k.
\end{equation}

\paragraph{Transformer Encoder.}
$L$ standard Transformer encoder layers are applied, each consisting of
multi-head self-attention (MHSA) and a feed-forward network (FFN):
\begin{align}
  \mathbf{A}^{(\ell)} &= \text{MHSA}\!\left(\mathbf{H}^{(\ell-1)}\right) + \mathbf{H}^{(\ell-1)}, \\
  \mathbf{H}^{(\ell)} &= \text{FFN}\!\left(\text{LN}(\mathbf{A}^{(\ell)})\right) + \text{LN}(\mathbf{A}^{(\ell)}),
\end{align}
where MHSA with $n_h$ heads computes:
\begin{equation}
  \text{MHSA}(\mathbf{H}) = \text{Concat}(\text{head}_1, \ldots, \text{head}_{n_h})\,\mathbf{W}^O,
  \quad
  \text{head}_j = \text{Softmax}\!\left(\frac{\mathbf{Q}_j\mathbf{K}_j^\top}{\sqrt{d_k}}\right)\mathbf{V}_j,
\end{equation}
with $\mathbf{Q}_j = \mathbf{H}\mathbf{W}_j^Q$, $\mathbf{K}_j = \mathbf{H}\mathbf{W}_j^K$,
$\mathbf{V}_j = \mathbf{H}\mathbf{W}_j^V$, and $d_k = d_h / n_h$.
The FFN uses GELU activation with a bottleneck factor of 2: $\text{FFN}(\mathbf{x}) = \mathbf{W}_2\,\text{GELU}(\mathbf{W}_1\,\mathbf{x})$.

\paragraph{Cross-Attention Pooling.}
A global query vector is formed by mean-pooling the encoder output, then used as
the query in a cross-attention operation:
\begin{equation}
  \mathbf{q} = \frac{1}{K}\sum_{k=1}^{K} \mathbf{H}^{(L)}_k,
  \quad
  \mathbf{z}_{\text{bag}}, \boldsymbol{\alpha} = \text{MHSA}(\mathbf{q}, \mathbf{H}^{(L)}, \mathbf{H}^{(L)}).
\end{equation}

% ------------------------------------------------------------
\subsection{Convolutional MIL}
\label{sec:conv}
% ------------------------------------------------------------

The convolutional variant treats the patch dimension as a spatial axis and
applies 1-D convolutions to model local patch neighbourhoods.

\paragraph{Patch Attention.}
Each patch embedding is scored by a 1-D convolutional attention layer.
The feature vector of patch $k$ is treated as a single-channel 1-D signal of
length $d$:
\begin{equation}
  s_k = \text{Conv1d}(\mathbf{x}_k) \in \mathbb{R},
\end{equation}
with masking of padded positions ($s_k \leftarrow -10^{10}$ for $m_k = 0$)
before softmax normalisation:
\begin{equation}
  \alpha_k = \frac{\exp(s_k)}{\sum_{j} \exp(s_j)}.
\end{equation}

\paragraph{Weighted Aggregation and Group Pooling.}
Attention-weighted patches $\hat{\mathbf{x}}_k = \alpha_k \mathbf{x}_k$ are
aggregated into $G$ ordered groups of approximately equal size:
\begin{equation}
  \mathbf{g}_g = \sum_{k \in \mathcal{G}_g} \hat{\mathbf{x}}_k,
  \quad g = 1, \ldots, G,
\end{equation}
where $\mathcal{G}_g = \{(g-1)\lceil K/G \rceil + 1, \ldots, \min(g\lceil K/G\rceil, K)\}$.
Alternatively, the top-$G$ patches ranked by $\alpha_k$ can be selected directly.

\paragraph{Convolutional Blocks.}
The grouped features $\mathbf{G} \in \mathbb{R}^{B \times G \times d}$ are
processed by a stack of residual 1-D ConvBlocks:
\begin{equation}
  \mathbf{G}^{(\ell+1)} = \text{ReLU}\!\left(\text{BN}\!\left(\text{Conv1d}(\mathbf{G}^{(\ell)})\right)\right) + \mathbf{W}_{\text{res}}\mathbf{G}^{(\ell)},
\end{equation}
followed by global average pooling over the group dimension and a two-layer
MLP classifier.

% ------------------------------------------------------------
\subsection{Dense MIL}
\label{sec:dense}
% ------------------------------------------------------------

The Dense MIL model uses the same patch attention and group aggregation as the
convolutional variant, but replaces the 1-D ConvBlocks with fully connected
DenseBlocks. The grouped feature matrix is flattened into a single vector:
\begin{equation}
  \mathbf{v} = \text{vec}(\mathbf{G}) \in \mathbb{R}^{G \cdot d},
\end{equation}
which is then passed through a configurable stack of residual DenseBlocks:
\begin{equation}
  \mathbf{v}^{(\ell+1)} = \sigma\!\left(\text{LN}\!\left(\mathbf{W}^{(\ell)}\mathbf{v}^{(\ell)}\right)\right) + \mathbf{W}_{\text{res}}^{(\ell)}\mathbf{v}^{(\ell)},
\end{equation}
where $\sigma$ is a configurable activation (ReLU, GELU, Leaky ReLU, or Tanh).
A final linear layer maps the last hidden state to class logits.

% ============================================================
\section{Training Procedure}
\label{sec:training}
% ============================================================

\subsection{Loss Function}

All models are trained to minimise the class-weighted cross-entropy loss:
\begin{equation}
  \mathcal{L} = -\frac{1}{N}\sum_{i=1}^{N} w_{y_i} \log p_{\theta}(y_i \mid B_i),
\end{equation}
where $p_{\theta}(y \mid B) = \text{Softmax}(\mathbf{z})_y$ and
$w_c = N / (C \cdot N_c)$ are inverse-frequency class weights with $N_c$ samples
in class $c$. When the dataset is only mildly imbalanced, equal weights
($w_c = 1$) may also be used.

\subsection{Optimiser}

The AdamW optimiser~\cite{loshchilov2019} is used with decoupled weight decay:
\begin{equation}
  \boldsymbol{\theta}_{t+1} = \boldsymbol{\theta}_t
    - \eta_t \left(\frac{\hat{\mathbf{m}}_t}{\sqrt{\hat{\mathbf{v}}_t} + \epsilon} + \lambda \boldsymbol{\theta}_t\right),
\end{equation}
where $\hat{\mathbf{m}}_t$ and $\hat{\mathbf{v}}_t$ are bias-corrected first and second
moment estimates of the gradient, $\eta_t$ is the current learning rate,
$\lambda$ is the weight decay coefficient, and $\epsilon$ is a numerical stability
constant. Default hyperparameters: $\eta_0 = 10^{-4}$, $\lambda = 10^{-4}$,
$\beta_1 = 0.9$, $\beta_2 = 0.999$.

\subsection{Learning Rate Scheduling}

A ReduceLROnPlateau scheduler monitors the validation loss and halves the
learning rate if no improvement is observed for 5 consecutive epochs:
\begin{equation}
  \eta_{t+1} = \begin{cases}
    \gamma \, \eta_t & \text{if no val. loss improvement for } p_{\text{sched}} \text{ epochs}, \\
    \eta_t           & \text{otherwise},
  \end{cases}
\end{equation}
with reduction factor $\gamma = 0.5$ and scheduler patience $p_{\text{sched}} = 5$.

\subsection{Early Stopping}

Training is halted if the primary selection metric does not improve for
$p_{\text{stop}}$ consecutive epochs (default $p_{\text{stop}} = 10$).
Two selection criteria are supported:
\begin{itemize}
  \item \textbf{F1 Macro} (primary): the macro-averaged F1 score over both classes.
        Higher is better. This criterion is preferred because it is class-balanced and
        more informative than accuracy under mild imbalance.
  \item \textbf{Validation loss}: standard cross-entropy on the held-out set.
        Lower is better.
\end{itemize}
The model state corresponding to the best observed metric value is restored
at the end of training.

\subsection{Regularisation}

Multiple regularisation mechanisms are used to prevent overfitting:
\begin{enumerate}
  \item \textbf{Dropout}: applied after projections, within the LSTM (inter-layer),
        inside the attention MLP, and before the classifier head.
  \item \textbf{Weight decay}: L2 penalty on all non-bias parameters via AdamW.
  \item \textbf{LayerNorm}: applied to LSTM outputs and intermediate classifier
        activations to stabilise training and reduce internal covariate shift.
  \item \textbf{Kaiming initialisation}: prevents vanishing/exploding gradients at
        the start of training.
\end{enumerate}

% ============================================================
\section{Evaluation Protocol}
\label{sec:eval}
% ============================================================

\subsection{Performance Metrics}

For each dataset split (train, validation, test), the following metrics are computed:
\begin{align}
  \text{Accuracy}  &= \frac{TP + TN}{TP + TN + FP + FN}, \\[4pt]
  \text{Precision} &= \frac{TP}{TP + FP}, \\[4pt]
  \text{Recall}    &= \frac{TP}{TP + FN}, \\[4pt]
  \text{F1}        &= \frac{2 \cdot \text{Precision} \cdot \text{Recall}}{\text{Precision} + \text{Recall}}, \\[4pt]
  \text{F1}_{\text{macro}} &= \frac{1}{C}\sum_{c=1}^{C} \text{F1}_c, \\[4pt]
  \text{AUC}       &= \int_0^1 \text{TPR}(t)\, \mathrm{d}\,\text{FPR}(t),
\end{align}
where $TP$, $TN$, $FP$, $FN$ denote true positives, true negatives, false positives,
and false negatives with respect to the positive (short-survivor) class.
The AUC is computed from the soft probability scores output by the model.

\subsection{Bootstrap Confidence Intervals}

All reported metrics are accompanied by empirical 95\% confidence intervals (CI)
computed via non-parametric bootstrap with $B = 1{,}000$ resamples. For a metric
$M$ computed on the test set of size $n$:
\begin{enumerate}
  \item Draw $B$ bootstrap samples $\mathcal{S}_b = \{i_1, \ldots, i_n\}$
        (with replacement) from $\{1, \ldots, n\}$.
  \item Compute $M^{(b)}$ on each $\mathcal{S}_b$.
  \item Report the $[\alpha/2, 1-\alpha/2]$ empirical percentiles of
        $\{M^{(b)}\}_{b=1}^{B}$ as the CI, with $\alpha = 0.05$.
\end{enumerate}

\subsection{Cross-Validation}

$k$-fold cross-validation (default $k = 5$) is supported for model selection and
variance estimation. The dataset is partitioned into $k$ stratified folds;
in each fold, one partition is held out for testing and the remainder is used for
training and validation. Metrics are aggregated across folds as mean $\pm$ standard
deviation.

\subsection{Centre-Based Subgroup Evaluation}

Because the dataset is collected from multiple clinical centres, performance is
additionally evaluated per centre. For each centre $c$ with at least
$n_{\min} = 10$ samples, all classification metrics are computed on the subset
$\mathcal{D}_c = \{(B_i, y_i) : \text{center}_i = c\}$. Centre-level metrics
are weighted by sample count to produce a grand average:
\begin{equation}
  \bar{M} = \frac{\sum_c |\mathcal{D}_c| \cdot M_c}{\sum_c |\mathcal{D}_c|}.
\end{equation}
This analysis exposes performance disparities driven by scanner heterogeneity,
acquisition protocol differences, and institutional case-mix variation.

% ============================================================
\section{Key Modelling Assumptions}
\label{sec:assumptions}
% ============================================================

We enumerate the primary assumptions underlying the framework and their
clinical/technical justifications.

\begin{enumerate}

  \item \textbf{Bag-level supervision is sufficient.}
        Patient-level binary labels (short vs.\ long survivor) are propagated to the
        entire bag without patch-level annotation. This is the fundamental MIL
        assumption: at least one (or collectively, the ensemble of) informative
        patch(es) is responsible for the bag label. For survival prediction, this is
        clinically plausible because the prognostically relevant signal (e.g., tumour
        heterogeneity, necrosis, vascular invasion) is distributed across multiple
        sub-volumes rather than localised to a single patch.

  \item \textbf{CT-FM embeddings are a sufficient statistic for patch-level risk.}
        The framework assumes that the 512-dimensional CT-FM embedding captures all
        radiologically prognostic information available at the patch level.
        Information not encoded by CT-FM (e.g., absolute Hounsfield values in the
        raw volume outside the patch, patient demographics) is not accessible to the
        MIL classifier. This motivates future work on multi-modal fusion.

  \item \textbf{Zero-padding does not introduce spurious signals.}
        Padded positions are represented as zero vectors $\mathbf{0} \in \mathbb{R}^d$.
        After the feature projection, the mask $\mathbf{m}$ is computed on the
        projected representation; any residual non-zero signal in padded positions
        (due to the projection bias $\mathbf{b}_{\text{proj}}$) is explicitly
        suppressed by mask multiplication. The attention softmax is also
        implicitly biased away from padded tokens because their LSTM outputs are
        zeroed before attention scoring.

  \item \textbf{Anatomical ordering is meaningful for the LSTM.}
        The superior-to-inferior ordering of patches is assumed to be the most
        clinically informative ordering for the LSTM, reflecting the continuous
        spatial extent of the tumour along the cranio-caudal axis. This is a
        \emph{strong inductive bias}: it implies that adjacent patches share
        anatomical context and that the temporal dependencies modelled by the
        LSTM correspond to genuine spatial dependencies in the CT volume.
        This assumption would break if patches were extracted non-contiguously or
        from multiple anatomical sites without a clear ordering.

  \item \textbf{Bidirectionality is beneficial.}
        Processing the patch sequence in both directions is assumed to be
        advantageous because prognostic signals may exist at any anatomical level,
        and the reverse pass allows each hidden state to be informed by patches
        that appear anatomically below it in the volume. The concatenation of
        forward and backward states doubles the effective hidden dimension.

  \item \textbf{Survival is a binary classification problem.}
        The continuous time-to-event nature of overall survival is reduced to a
        binary threshold problem. This discards censoring information and ignores
        the time dimension beyond the threshold. For clinical deployment, a
        survival regression or hazard estimation approach may be preferable.

  \item \textbf{Multi-centre data follows a common feature distribution.}
        Although data originate from different clinical centres (with potentially
        different scanners, acquisition protocols, and reconstruction kernels),
        CT-FM is assumed to produce embeddings that are sufficiently invariant to
        these nuisances. Centre-based evaluation (\cref{sec:eval}) is performed
        specifically to audit this assumption.

\end{enumerate}

% ============================================================
\section{Hyperparameter Summary}
\label{sec:hyperparams}
% ============================================================

\Cref{tab:hyperparams} summarises the default hyperparameters of the LSTM-based
MIL model and the training procedure.

\begin{table}[h]
\centering
\caption{Default hyperparameters for the LSTM-based MIL model and training.}
\label{tab:hyperparams}
\begin{tabular}{lll}
\toprule
\textbf{Parameter} & \textbf{Symbol} & \textbf{Default value} \\
\midrule
\multicolumn{3}{l}{\textit{Embedding}} \\
Input embedding dimension     & $d$               & 512 \\
Maximum patches per bag       & $K_{\max}$        & 300 \\
\midrule
\multicolumn{3}{l}{\textit{LSTM}} \\
Hidden dimension              & $d_h$             & 128 \\
Number of LSTM layers         & $L$               & 2 \\
Bidirectional                 & ---               & Yes ($2d_h = 256$ output) \\
Dropout (inter-layer)         & $p$               & 0.3 \\
\midrule
\multicolumn{3}{l}{\textit{Attention}} \\
Attention hidden dimension    & $d_h / 2$         & 64 \\
\midrule
\multicolumn{3}{l}{\textit{Classifier}} \\
Intermediate dimension        & $d_h$             & 128 \\
Output classes                & $C$               & 2 \\
\midrule
\multicolumn{3}{l}{\textit{Training}} \\
Optimiser                     & ---               & AdamW \\
Learning rate                 & $\eta_0$          & $10^{-4}$ \\
Weight decay                  & $\lambda$         & $10^{-4}$ \\
LR reduction factor           & $\gamma$          & 0.5 \\
Scheduler patience            & $p_{\text{sched}}$& 5 epochs \\
Early stopping patience       & $p_{\text{stop}}$ & 10 epochs \\
Maximum epochs                & $T_{\max}$        & 100 \\
Batch size                    & $B$               & 32 \\
Selection metric              & ---               & F1 Macro \\
\midrule
\multicolumn{3}{l}{\textit{Evaluation}} \\
Bootstrap resamples           & $B_{\text{boot}}$ & 1{,}000 \\
Confidence level              & ---               & 95\% \\
Cross-validation folds        & $k$               & 5 \\
\bottomrule
\end{tabular}
\end{table}

% ============================================================
\section{Discussion}
\label{sec:discussion}
% ============================================================

\subsection{The LSTM as a Spatial Risk Integrator}

The core conceptual contribution of the LSTM-based MIL model is its explicit
treatment of the CT volume as an ordered sequence of anatomical observations.
Whereas attention-based MIL models~\cite{ilse2018} and transformer-based
models treat the patch bag as an unordered set and learn permutation-invariant
aggregations, the LSTM model imposes a spatial prior: lesion-relevant information
accumulates progressively as the model traverses the volume from apex to base.

The cell state $\mathbf{c}_k$ at step $k$ can be interpreted as the model's
\emph{running estimate of cumulative risk} given all patches seen so far.
The forget gate $\mathbf{f}_k$ learns to retain this estimate across anatomically
benign regions, while the input gate $\mathbf{i}_k$ modulates the update magnitude
in response to locally abnormal patches. The final attention-pooled representation
$\mathbf{z}_{\text{bag}}$ thus synthesises spatial context from the entire volume
in a way that respects the anatomical ordering of the patches.

This architecture is particularly well-suited to predicting overall survival
in lung cancer, where the prognostic signal arises from the interplay of
tumour burden across multiple lobes, mediastinal involvement, and the regional
extent of disease --- all of which are spatially distributed phenomena.

\subsection{Comparison of Aggregation Strategies}

The framework includes four complementary aggregation strategies:
\begin{itemize}
  \item \textbf{LSTM}: exploits sequential structure; captures long-range
        anatomical dependencies; $\mathcal{O}(K \cdot d_h)$ memory per sample.
  \item \textbf{Transformer}: permutation-invariant; captures all pairwise
        patch interactions; $\mathcal{O}(K^2)$ self-attention cost.
  \item \textbf{Convolutional}: captures local patch-neighbourhood structure;
        efficient; limited long-range modelling capacity.
  \item \textbf{Dense}: no spatial structure assumed beyond group aggregation;
        simplest inductive bias; highest risk of overfitting on small datasets.
\end{itemize}

\subsection{Limitations}

\begin{itemize}
  \item \textbf{Censoring}: binary thresholding of OS discards censored observations
        and the temporal dimension of survival data.
  \item \textbf{Fixed encoder}: the CT-FM backbone is frozen; joint fine-tuning
        may improve task-specific feature alignment at the cost of increased
        compute and overfitting risk on small cohorts.
  \item \textbf{Patch ordering sensitivity}: the LSTM's performance depends on a
        consistent and clinically meaningful patch ordering; inconsistencies in
        patch extraction across centres could degrade performance.
  \item \textbf{No explicit anatomical localisation}: the framework does not
        enforce tumour-centred patch extraction; the model must learn to attend to
        relevant regions from weakly supervised signal alone.
\end{itemize}

% ============================================================
\section*{References}
% ============================================================

\begin{thebibliography}{9}

\bibitem{carbonneau2018}
M.-A.\ Carbonneau, V.\ Cheplygina, E.\ Granger, and G.\ Gagnon,
``Multiple instance learning: A survey of problem characteristics and applications,''
\textit{Pattern Recognition}, vol.\ 77, pp.\ 329--353, 2018.

\bibitem{ilse2018}
M.\ Ilse, J.\ M.\ Tomczak, and M.\ Welling,
``Attention-based deep multiple instance learning,''
in \textit{Proc.\ ICML}, 2018.

\bibitem{he2015}
K.\ He, X.\ Zhang, S.\ Ren, and J.\ Sun,
``Delving deep into rectifiers: Surpassing human-level performance on ImageNet classification,''
in \textit{Proc.\ ICCV}, 2015.

\bibitem{loshchilov2019}
I.\ Loshchilov and F.\ Hutter,
``Decoupled weight decay regularization,''
in \textit{Proc.\ ICLR}, 2019.

\bibitem{ctfm}
CT-FM: A Foundation Model for CT Imaging,
pretrained on $>$140{,}000 CT scans from all anatomical sites.

\end{thebibliography}

\end{document}
